\begin{abstract}
당구에서 공을 넣는 샷은 연속적인 action 공간 안의 아주 좁은 영역에만 존재하며, 조준이 조금만 어긋나도 결과가 완전히 달라진다.
이 때문에 기존 접근은 각자 한계에 부딪힌다.
표준 강화학습은 action 공간 전체에 대해 하나의 Q-function을 학습해야 해서 이 좁은 영역을 정확히 짚어내지 못하고,
search 기반 planner는 강하지만 샷마다 많은 계산이 필요하며,
기하 공식 기반 방법은 사람의 당구 지식에 의존한다.

우리는 사람의 지식 없이 학습하면서 search 없이 바로 샷을 고르는 당구 정책을 제안한다.
모델은 공과 포켓의 각 조합을 GNN으로 표현하고, 조합마다 겨눌 방향을 제안한 뒤 그 샷이 들어갈지를 예측해 가장 좋은 샷을 고른다.
하나의 Q로 action 공간 전체를 다루는 대신, 문제를 ``어디를 겨눌까''와 ``들어갈까''로 나눈 것이다.
학습에는 teacher나 기하 공식 없이 게임 결과만 쓰며, 모델은 자기가 둔 샷의 결과로 다시 학습하면서 스스로 개선된다.

이렇게 학습한 정책은 network forward 한 번으로 샷을 결정한다.
3공을 모두 넣는 데 평균 3.20샷이 걸려 search 기반 planner(2.98샷)에 근접하고 표준 강화학습(7.89샷)을 크게 앞선다.
또 공--포켓 조합 단위로 학습했기 때문에, 3공에서만 학습한 모델이 6공과 9공 테이블에서도 그대로 동작하며,
이때 run-out 비율은 공식 기반 모델보다 높다.
\end{abstract}
